% Ad M. Brutum 2.1 (ad-brutum-02-01)
% Source: https://raw.githubusercontent.com/PerseusDL/canonical-latinLit/master/data/phi0474/phi059/phi0474.phi059.perseus-lat1.xml
% Fetched by scripts/fetch_latin.py

% Dateline (Perseus): Scr. Romae ex. m. Mart. ant in. Apr. a. 711 (43) .

\ciceroLetterOpener{Cicero Bruto salutem}

\ciceroSection{1} Cum haec scribebam, res existimabatur in extremum adducta discrimen. tristes enim de Bruto nostro litterae nuntiique adferebantur. me quidem non maxime conturbabant. his enim exercitibus ducibusque quos habemus nullo modo poteram diffidere neque adsentiebar maiori parti hominum. fidem enim consulum non condemnabam quae suspecta vehementer erat; desiderabam non nullis in rebus prudentiam et celeritatem; qua si essent usi, iam pridem rem publicam reciperassemus. non enim ignoras quanta momenta sint in re publica temporum et quid intersit idem illud utrum ante an post decernatur, suscipiatur, agatur. omnia quae severe decreta sunt hoc tumultu, si aut quo die dixi sententiam perfecta essent et non in diem ex die dilata aut, quo ex tempore suscepta sunt ut agerentur, non tardata et procrastinata, bellum iam nullum haberemus.

\ciceroSection{2} omnia , Brute, praestiti rei publicae quae praestare debuit is qui esset in eo in quo ego sum gradu senatus populique iudicio conlocatus, nec illa modo quae nimirum sola ab homine sunt postulanda, fidem, vigilantiam, patriae caritatem. ea sunt enim quae nemo est qui non praestare debeat. ego autem ei qui sententiam dicat in principibus de re publica puto etiam prudentiam esse praestandam nec me, cum mihi tantum sumpserim ut gubernacula rei publicae prehenderem, minus putarim reprehendendum si inutiliter aliquid senatui suaserim quam si infideliter.

\ciceroSection{3} Acta quae sint quaeque agantur scio perscribi ad te diligenter; ex me autem illud est quod te velim habere cognitum, meum quidem animum in acie esse neque respectum ullum quaerere nisi me utilitas civitatis forte converterit; maioris autem partis animi te Cassiumque respiciunt. quam ob rem ita te para, Brute, ut intellegas aut, si hoc tempore bene res gesta sit, tibi meliorem rem publicam esse faciendam aut, si quid offensum sit, per te esse eandem reciperandam.
